\chapter{Integración y aproximación}
\label{cap:integracion-aproximacion}

La integral de Riemann se construye mediante particiones y sumas superiores e inferiores. El capítulo reúne los criterios de integrabilidad, el teorema fundamental del cálculo y las aproximaciones locales de Taylor con control explícito del resto \cite{apostol1974,spivak2008,rudin1976}.

\section{Particiones de un intervalo}
Particiones, refinamientos, malla y variación de una función en subintervalos.

\section{Sumas superiores e inferiores}
Sumas de Darboux y criterio de integrabilidad.

\section{Integral de Riemann}
Definición, existencia para funciones acotadas y ejemplos básicos.

\section{Criterios de integrabilidad}
Integrabilidad de funciones continuas y monótonas y criterio de Lebesgue para Riemann.

\section{Propiedades de la integral}
Linealidad, monotonía, aditividad por intervalos y estimaciones.

\section{Teorema fundamental del cálculo}
Relación entre derivación e integración y sus dos partes.

\section{Fórmula de Taylor}
Polinomios de Taylor y expansión alrededor de un punto.

\section{Teorema de Taylor}
Hipótesis de regularidad y representación del error.

\section{Formas del resto}
Resto de Lagrange, forma integral y otras expresiones del residuo.

\section{Estimaciones del error}
Cotas del resto y aplicaciones a aproximaciones cuantitativas.
